\section{Discussion and limitations}
\label{sec:disc}

\subsection{The C++ question}
\label{sec:disc:cpp}

\sysname is written in C++, not Rust, and we would rather address
that here than in the rebuttal. The pitch is not that C++ is
fine---we hold no illusions about its memory-safety record---but that
capability discipline can be retrofitted onto an unsafe-language
ecosystem, a direction the security literature mostly leaves
unexplored. \sysname is evidence that a 4096-token capability layer
with constant-time mediation can be built over a C++ runtime without
the language becoming the obstacle. That is a different bet from the
one Tock~\cite{tock}, Theseus~\cite{theseus}, and Redox make, all of
which decide that the language is the point and choose Rust.

The bet is testable, and we test it: the CBMC proof, the
property-test oracle, the libFuzzer harness, and the constant-time
NASM primitives together cover the surface where C++'s unsafety most
plausibly turns into a security bug. The residual risk is stated
plainly---an attacker who corrupts the runtime's own memory owns
every token---and the paper's contribution is scoped accordingly: the
upper-half capability story, not the lower-half language-safety one.

A Rust port of the hot path is on our roadmap, contingent on
reviewers telling us the language is the paper's primary weakness.

\subsection{The Aeron gap}
\label{sec:disc:aeron}

Aeron's published 250\,ns RTT is the tightest single-host IPC number
in production use. Our gap of \aeronGapPercent\% is the figure we
most want to defend downward, and the one reviewers will push
upward. It decomposes into three parts:

\begin{itemize}
\itemsep0pt
\item Two \texttt{validate()} calls per round trip
($\sim$\gateOverheadTail\,ns each at p99). Irreducible without
dropping security or weakening the constant-time comparator.
\item C++23 range-checked containers and \texttt{Result<T>} unwraps
on the hot path. Aeron's published number comes from a leaner,
C-flavoured API; a no-stdlib rework of our hot path would claw some
of this back, and we have not done it.
\item Fence policy. Aeron's ring uses release-acquire on
\texttt{write\_index} alone; our two-phase MPSC claim/commit needs it
on \texttt{write\_claim\_index} as well. That is a correctness
requirement, and it costs cycles relative to Aeron's SPSC-only
arrangement.
\end{itemize}

We do not expect to beat Aeron at raw RTT---the design template is
shared---but the gap should narrow with a hot-path API streamline.

\subsection{Reading HongMeng}
\label{sec:disc:hongmeng}

HongMeng (OSDI~'24) is the nearest competitor on the capability +
Linux + fast-IPC triple, and the honest framing is that it owns the
kernel-level cell while \sysname owns the userspace one. HongMeng's
authors describe walking back chain-revocable capabilities to
``address tokens'' because revocation cost hurt their common case. We
read that as evidence about \emph{their point in the design
space}---in-kernel, full revocation---being expensive, not as
evidence that revocation is unaffordable in general. Sitting in
userspace, \sysname pays a different price (the slot load on every
message) and in exchange keeps full chain revocation with $O(1)$
time-to-effect.

\subsection{Single-machine scope}
\label{sec:disc:scope}

Everything here is single-host. The cross-host
story---attestation-backed delegation, remote token transfer, tokens
as attestation artefacts inside TEEs---belongs to follow-up work
(M-23 ConfCompute), excluded deliberately to keep this contribution
sharp. Readers who want the TEE-binding design will find it in the
architecture overview~\cite{astra-confcompute}, not in this paper's
measurements.

\subsection{No macrobenchmark}
\label{sec:disc:macro}

We ship no NGINX/Redis/SQLite macrobenchmark, and the reason is
space: five transports head-to-head, the $O(1)$ confirmation,
revocation latency, and perf counters already fill the twelve pages.
A macrobenchmark would also shift the headline from ``the gate is
cheap'' to ``the gate is cheap for this particular application,''
which is the weaker claim. The harness scaffolding ships with the
artefact, and the macrobenchmark numbers will appear in the
artefact-evaluation appendix.
